\documentclass[11pt]{article}
\usepackage{icmlko}
\begin{document}

\title{열여섯 행은 자기에게 해롭고 남에게 이롭다 --- 사전 등록이 뒤집혔는데도 기각한 까닭}
\author{월드모델 실험실}
\date{노트 578 · 2026년 8월 5일}
\maketitle

\begin{abstract}
학습 문턱 $15 \to 40$ 은 이 판에서 \textbf{팝업 하나에만 닿는다}(학습
$16$ 행, 다음으로 적은 아이돌이 $54$). 노트 $550$ 이 사전 등록했다 ---
``빼면 팝업 $\rho$ 가 내려간다, $16$ 행이라도 배우는 게 낫다''.

\textbf{반대였다.}

\begin{center}
\begin{tabular}{lrrl}
\hline
자 & 차(문턱 $40 -$ $15$) & 부호 & \\
\hline
\textbf{팝업} & $\mathbf{+0.0461}$ & $\mathbf{12/12}$ & \textbf{0 밖 양수} \\
비게임 앱 & $+0.0052$ & $10/12$ & 0 밖 양수 \\
판 & $-0.0004$ & $6/12$ & 0을 문다 \\
날짜 통제 판 & $-0.0005$ & $8/12$ & 0을 문다 \\
KR 만화 & $-0.0029$ & $5/12$ & 0을 문다 \\
\textbf{무리 안 판} & $\mathbf{-0.0139}$ & $\mathbf{0/12}$ & \textbf{0 밖 음수} \\
\hline
\end{tabular}
\end{center}

$16$ 행으로 배운 것이 \textbf{ALL5 집 밖 예보보다 나쁘다.} 그런데
\textbf{다섯 자 중 무리 안 판이 $0$ 밖 음수}라 기각했다 --- \textbf{문턱
$15$ 유지.}
\end{abstract}

\begin{figure}[t]
\centering
\includegraphics[width=\textwidth]{figs/f1.pdf}
\caption{\textbf{가} --- 팝업만 오르고 무리 안이 내린다.
\textbf{나} --- 장부. \textbf{다} --- 뒤집혔는데도 기각.}
\end{figure}

\section{장부를 맞춰 보면}

팝업 유보는 $65$ 행, 분모 $3{,}369$ 의 \textbf{$1.9\%$} 다. 팝업이
$+0.0461$ 오르면 판에는 $0.0461 \times 0.0193 = \mathbf{+0.0009}$ 가
들어와야 한다.

\begin{center}
\begin{tabular}{lr}
\hline
 & 판 $\rho$ 몫 \\
\hline
팝업이 내놓는 것 & $+0.0009$ \\
판에 실제로 남은 것 & $\mathbf{-0.0004}$ \\
\hline
\textbf{나머지 열이 가져간 것} & $\mathbf{+0.0013}$ \\
\hline
\end{tabular}
\end{center}

\begin{quote}
\textbf{열여섯 행짜리 도메인은 순 기부자다.} 자기를 적합하기엔 모자라서
\emph{자기에게는 해로운데}, 그 열여섯 행이 공유 축의 눈금을 맞춰
\emph{남에게는 이롭다}. 무리 안 판 $-0.0139$($0/12$, 이 세션에서 제일
선명한 부호)가 그 자리를 가리킨다.
\end{quote}

도메인별로 보면 팝업을 빼서 손해 보는 쪽이 \textbf{도서}($-0.0145$)와
\textbf{아이돌}($-0.0132$)이다. 둘 다 유보가 작다($163$ $\cdot$ $51$).
\textbf{작은 것끼리 서로 받쳐 주고 있다.}

\section{왜 뒤집혔는데도 기각인가}

사전 등록의 (1)이 반대로 나왔으니 규칙대로면 ``$40$ 으로 올린다''였다.
올리지 않았다. \textbf{(3)이 걸렸기 때문이다} --- 무리 안 판이 $0/12$ 로
음수다.

\begin{quote}
\textbf{사전 등록은 하나가 아니라 셋이었고, 셋을 다 지킨다.} 뒤집힌
가설 하나가 나머지 둘을 무효로 만들지 않는다. 이 판이 관문을 \emph{다섯
자}로 둔 이유가 그것이다 --- \textbf{한 자가 크게 움직였다고 채택하지
않는다}(노트 $518$).
\end{quote}

거울(노트 $426$)도 안 섰다 --- KR 만화가 $0$ 을 물고($5/12$) 비게임 앱만
양수다. \textbf{집 밖 짝 둘이 같은 말을 하지 않는다.}

\section*{갈래 하나를 닫는다}

``학습 문턱 $40 \to 15$ 를 전 판에서 재확정''은 이 판의 상시 갈래였다.
닫는다 --- 판 $-0.0004$ $[-0.0022,+0.0015]$ 로 \textbf{전 판에서
중립}이고, 다른 자가 반대하므로 \textbf{$15$ 를 유지한다.} 챔피언
\texttt{MIN\_TRAIN=15} 는 그대로다.

\section*{조항}

\begin{quote}
\textbf{도메인의 값어치는 그 도메인 점수로 재지 않는다.} 팝업은 자기
점수를 $+0.046$ 올리는 대신 판에서 $-0.0013$ 을 뺏는다. \textbf{묶어
배우는 판에서 ``이 도메인이 도움이 되나''는 그 도메인을 빼 보고 \emph{남}을
보는 물음이다.}
\end{quote}

\section*{남은 것}

\begin{itemize}
\item \textbf{팝업만 빼되 그 축은 남기는 팔} --- 지금 팔은 통째로 뺀다.
$16$ 행이 주는 것이 \emph{라벨}인지 \emph{축의 눈금}인지 갈린다.
\item \textbf{무리 안 판이 왜 그렇게 민감한가} --- $0/12$ 는 이 세션에서
제일 선명한 부호다. 무리 안 판만 크게 움직이는 팔은 드물다.
\item \textbf{팝업과 도서} --- 도서가 $-0.0145$ 로 제일 크게 손해 본다.
둘이 무슨 축을 나눠 쓰는지 안 봤다.
\end{itemize}

\end{document}
